\documentclass[12pt]{article}

% set margins and spacing
\addtolength{\textwidth}{1.3in}
\addtolength{\oddsidemargin}{-.65in} %left margin
\addtolength{\evensidemargin}{-.65in}
\setlength{\textheight}{9in}
\setlength{\topmargin}{-.5in}
\setlength{\headheight}{0.0in}
\setlength{\footskip}{.375in}
\renewcommand{\baselinestretch}{1.0}
\linespread{1.0}

% load miscellaneous packages
\usepackage{csquotes}
\usepackage[american]{babel}
\usepackage[usenames,dvipsnames]{color}
\usepackage{graphicx,amsbsy,amssymb, amsmath, amsthm, MnSymbol,bbding,times, verbatim,bm,pifont,pdfsync,setspace,natbib}
\usepackage{multirow, amsmath}
\usepackage[table]{xcolor}
\usepackage{float}

% enable hyperlinks and table of contents
\usepackage[pdftex,
bookmarks=true,
bookmarksnumbered=false,
pdfview=fitH,
bookmarksopen=true,hyperfootnotes=false]{hyperref}
\hypersetup{colorlinks=true, linkcolor=blue, citecolor=blue}

% enable flowchart diagram package
\usepackage{tikz}
\usetikzlibrary{shapes.geometric, arrows}
\usepackage{graphicx}

% define environments
\newtheorem{definition}{Definition}
\newtheorem{fact}{Fact}
\newtheorem{result}{Result}
\newtheorem{proposition}{Proposition}



\begin{document}
\title{Crime and Drug Treatment across Detroit}
\author{Rachel Gaudreau\thanks{Syracuse University, Economics Department. Email: regaudre@syr.edu.} \and Sophia Oritz-Heaney\thanks{Syracuse University, Economics Department. Email: sgortizh@syr.edu} \and Weston Maechling\thanks{Syracuse University, Economics Department. Email: wrmaechl@syr.edu.} \and Leo Gershman\thanks{abc}}
\date{\vskip-.1in \today}
\maketitle

\bigskip
\begin{center} {\bf Abstract} \end{center}

\begin{quote}
    {\small We examine the effects drug and alcohol treatment centers have on their neighbors in Detroit, Michigan. In this study, we find the relationship between the distance from a substance abuse treatment centers (SATC) and the number of 911 calls in an area. This paper creates a framework that locates all the  substance abuse treatment centers in the greater Detroit area and then creates rings around them at distances ranging from 50 to 2500 meters to find the number of calls within each ring. Our analysis uses multiple t tests to find the number of calls within a ring and the probability that this is indicative of all crime in the city throughout its history. This data was found through the "City of Detroit's Open Data Portal". Our group looked in the "Police Serviced 911 Calls" section. This portal contains a wealth of information on the type of crime committed along with where and when it was committed along with the type of crime. With this information, our group finds the correlation between the distance to a SATC and 911 call. With this information, a t test is conducted. We did this for all 911 calls and for calls that involved crimes. The results of this experiment were different than our hypothesis. These results show that that there are more 911 calls and crimes reported the closer it is to a SATC. This could be because SATCs are placed in areas where they can help the most. Our group believes that they are placed in worse neighborhoods that have more 911 calls and report more crimes. Along with this, due to "Not in my backyard" (NIMBY), a characterization of opposition by residents to infrastructure development in their local area, SATCs are only built in poorer neighborhoods. Information on the trend of 911 calls in relation to SATCs could help infrastructure planning in similar urban areas.} 
\end{quote}

\bigskip
\section{Introduction} \label{sec:introduction}
A common concern among the general American public is that crime rates are increasing. About 6 in 10 Americans share this concern according to Pew Research Center. Additionally, drug usage is another public concern in the American populous, especially in urban centers. This paper seeks to find a correlation between proximity to  substance abuse treatment centers and crime. Our hypothesis theorizes that a closer proximity to substance abuse treatment centers increases access to drug rehabilitation and mental health resources, which, in turn, decreases drug usage. This reduction in drug usage would theoretically lead to a decrease in drug-related crimes and, consequently, fewer calls to 911. So, we explored the question, "How does the distance to the closest drug rehabilitation center relate to the number of calls to 911?" 

Our hypothesis is based on the established impact of drug rehabilitation centers on community well-being . \cite{SAT_centers_and_crime} agrees with this point by showing that there is a decline on all violent crimes within the year of a substance abuse treatment center opening in an area.  Studies such as \cite{drugs_and_crime} and \cite{drugs_crime_space_time} have demonstrated that these centers can lower crime rates and elevate the standard of living in surrounding areas. This means that substance abuse treatment centers can help community well-being by doing more than preventing crime. This is shown in \cite{mental_healthcare_and_crime}, an article showing that drug dependency is a leading cause of criminal behavior. By reducing this dependency, substance abuse treatment centers improve the crime rates in their neighborhoods. The lack of criminal behavior allows people to fully embrace better lives that are closer to their communities further decreasing the amount of 911 calls. 

To explore the relationship between crime and SATC, we conducted a spatial analysis of Wayne County, where Detroit is located. We mapped the locations of drug rehabilitation centers in conjunction with the data points representing 911 calls. Our results showed that areas closer to substance abuse treatment centers had a higher number of 911 calls when compared to areas further away.  Interestingly, further analysis revealed that this trend persisted even when isolating calls related to criminal activities. This data does not support our hypothesis. Nonetheless, as argued in \cite{Socioeconomic-Determinants}, a higher amount of drug usage in an area causes an increase in drug usage. Building off of this, areas of higher crime, and thus drug usage, are closer to SATCs instead of further away from them as previously argued. 

A potential explanation for these findings could be the reluctance of wealthier communities to give drug rehabilitation centers consent inside their neighborhoods. They could fear that such facilities might attract individuals they perceive as a threat due to their unstable nature as addicts as shown in \cite{mental_health_and_disability}.  Along with this, these locations are selected to be in a place where they could have the most immediate impact. These are places that had the highest rates of crime and drug usage before the treatment centers were built. SATCs were built in areas that already had prevalent drug issues. These issues don't go away immediately.  From that point of view, it is completely logical that the SATCs would be near the places with the highest drug usage. This is where they are the most effective. 

Answering this question is crucial for developing strategies to mitigate drug-related crimes along with crimes as a whole in American cities. As shown in \cite{SAT_centers_and_crime}, drugs are a major source of civil unrest, social discord, and crime, particularly in urban areas. This is especially true in areas that, historically, have not been supported by government resources. By understanding the relationship between treatment center proximity and emergency call frequency, we can work toward solutions that foster safer and healthier communities. Our paper, attempts to provide further insight into the relationship between SATCs and crime in Detroit. In doing so, we hope to further the field of social economics and the consequences of SATC construction in an area. Although more can be done on the subject by adding more cities or seeing the relationship between the amount of 911 calls over the years, our study shows how SATCs are commonly found in areas that have a 911 calls along with explanations on why that could be the case.



\section{Literature Review} \label{sec:literature}
The geographical relationship between crime and access to mental health care is a widely studied topic but specific studies are widespread and disconnected. For example, mental illness is associated with an increased likelihood of substance abuse, which directly and indirectly relates an increase in crime to mental illness. One viable method of treating mental illness is through office-based mental healthcare. Thus, evidence shows that increasing access to office-based mental health care by 10 mental health care provider offices in a county can lead to a 0.4\% reduction in the county crime rate. \cite{mental_healthcare_and_crime}). An increase in access to mental health care also shows a connection to participation in government disability programs. This is largely because better access to mental health facilities increases the likelihood of proper diagnosis of serious mental illness which leads to access to programs such as Supplemental Security Income (SSI) and Social Security Disability Insurance (SS(D)I). \cite{mental_health_and_disability}. While access to disability insurance does not directly relate to crime this shows that increasing access to mental health care allows for more opportunities for support and better access to the social safety net. Furthermore, this effect is also evident when specifically considering access to substance abuse treatment (SAT) facilities. Evidence shows that increasing access to SAT facilities reduces violent crime due to the effectiveness in reducing crimes motivated by obtaining money to purchase drugs, reducing violence among those in the drug trade, and reducing drug usage which eases aggressive behavior due to drug abuse \cite{SAT_centers_and_crime}. 

However, a significant barrier to treatment could be the price of substance abuse treatment. In a study conducted between 2002 and 2003, the median willingness to pay for drug rehabilitation is shown to be below the average cost of treatment. Moreover, the cost of drug treatment is a significant indicator of self-reported probability of enrollment. \cite{cost_of_drug_treatment}. While it is likely that cost of treatment has changed in the subsequent decades, it is still important to acknowledge barriers to treatment when considering its affects on crime rates in an urban area, such as Detroit. 

Barriers to  treatment also contribute to evidence of a positive feedback loop between drug usage, crime, and standard of living in a given area. The drug trade, specifically the sale and use of drugs in a given area has been closely linked to an increase in the crime rate in the subsequent area. A case study of Miami-Dade county shows that drug activity on a given block can alter or nullify the relationship between household income and crime. Thus, drug usage can be a predictor of crime rates in a geographical area \cite{drugs_crime_space_time}. This effect also has implications on an areas living standards. Since, higher levels of drug usage leads to higher crime rates in a given geographical area, this corresponds to lower standards of living. This effect is due to the addictive properties of drugs, the numbing of decision making caused by drug usage, and the relationship between areas with lower standards of living and a higher tendency for those who live within such an area to turn to drugs as a means for escape \cite{drugs_and_crime}.

Considering the almost conflicting theories on how drugs, crime, and drug treatment are related and the affects on geographical area, understanding how drug treatment relates directly to crime in the area will give insight into possible policies to reduce crime in an area. Furthermore, given the lack of economic research on this topic we decided to focus on a single case study, in this case, Detroit, and the immediate surrounding communities to analyze the effects drug treatment centers have on crime their surrounding area. Moreover, given our focus on crime it is important to understand the impact the COVID-19 Pandemic had on crime in general. Practices such as Social Distancing and at-home quarantines led people to spend more time in privately occupied spaces. This led to a decrease in the amount of crimes committed in a public setting, such as robbery, homicide, physical violence, etc. However, this also led to an increase in crimes committed in private settings, such as domestic violence \cite{covid_and_crime}. Since, COVID-19 had a varied affect on crime for the purposes of our analysis we decided to limit our study of crime prior to the COVID-19 Pandemic.

\section{Theoretical Analysis}
\label{sec:theory}
This paper assesses the hypothesis that as you get further away from a substance abuse treatment center in Detroit, MI you can expect a higher number of calls to 911. \cite{SAT_centers_and_crime} demonstrate how crime rate is negatively related to the number of substance abuse treatment centers. As more centers open, they can provide treatment to people who were previously unable to receive any due to distance constraints. The paper also showed that this number was found through the likelihood dependency on emergency services decreasing, not just crime. This follows along with what \cite{Socioeconomic-Determinants} found;  lower drug usage in an area causes the people who live there to have more stable and prosperous lives. It shows that as someone frees themselves from drug addiction and increases their socioeconomic standing, they become more productive members of society. This greater access to treatment caused overall crime rates to drop as there was a reduction in crimes motivated by drugs or getting the money to acquire more. We can connect it with the findings of \cite{drugs_and_crime} who found that a decreased drug use in an area has a positive correlation with the change in crime in the same area. This rise in both drug usage and crime causes a negatively related change in the standards of living of that area which they found will again negatively relate to the amount of drug usage restarting the cycle. This framework allowed us to analyze our claim (see Figure 1) that the distance to the closest treatment center will impact the ability to receive treatment, which will decrease the level of drug use in that area, leading to a positively related effect on drug-related crimes and also positively effecting the total number of 911. 



\section{Data}
\label{sec:data}

\subsection{911 Calls in Detroit Area}

This data comes from the City of Detroit Open Data Portal's   \href{https://data.detroitmi.gov/datasets/detroitmi::police-serviced-911-calls/about}{Police Serviced 911 Calls}. It is generated by the Detroit Police Department's Crime Data Analytics when a call is placed. This data set covers 911 calls that are received by precincts around the Detroit metropolitan area from the year 2016-2022. There are 5,238,111 observations. Key variables include incident address, year, latitude, longitude, call category and call description. 

We isolated the observations to one year, 2017, from the original dataset so we could perform cross-sectional analysis. We chose observations from 2017 since the year a) overlaps with the second set of data and b) is before COVID-19, where crime rates statistics shifted in general \cite{covid_and_crime}. This is to eliminate COVID-19 effects from affecting our data analysis. Years 2016, 2018 and 2019 could also work for cross sectional analysis, but for simplicity, we only analyze one year in this report. The data includes latitude and longitude variables which we used to geocode each call to a single point in space.
\subsection{Substance Abuse Treatment Centers (SATC) in Detroit Area (Dataset 2)}

This data is from the Substance Abuse and Mental Health Services Administration (SAMHSA). Data is generated through the National Survey of Substance Abuse Treatment Services \href{https://www.samhsa.gov/data/data-we-collect/n-ssats-national-survey-substance-abuse-treatment-services}{N-SSATS} an annual survey of all known public and private substance abuse treatment services in the United States. SAMHSA conducts this survey annually. This data covers the names and addresses of substance abuse treatment services in operation in the Detroit Metropolitan area between the years 2015 to 2021. 

 
We used layered the datasets and geocode their locations. Using the coordinates (longitude and latitude variables) of each point, we created two new variables. Every 911 call is assigned one SATC, the one geographically closest to it. This means no observation is counted twice. 

We create distance rings, referred to as distance groups, and perform our analysis. We normalize the data to account for differences in total area of a distance group might alter results. We did this by dividing the area in squared meters by the total count of 911 calls in each distance group.

\subsection{Measuring Distance of each 911 Call to each SATC}
 \begin{tikzpicture}
% Define the radii corresponding to the first 3 distances (in meters)

\def\radiusOne{0.5}
\def\radiusTwo{1}
\def\radiusThree{1.75}


% Draw the concentric circles
\foreach \r/\label in {\radiusOne/, 
    \radiusTwo/,
    \radiusThree/,}
 {
\draw[thick] (0,0) circle (\r);
    \node at (\r, 0) {\label};
    }
\end{tikzpicture}

We created concentric rings around each SATC. Each 911 point falls within one circle, referred to as a distance ring or distance group, and is counted towards the SATC in its circle's center.

%Describe your data. Where you got it from, how it was generated, what variables you'll use, what data cleaning steps you had to take, where your processed data, code and documentation is stored.
%In a published paper, a lot of this detail will be in a data appendix. For the purposes of this report, include it all here (this may be the longest section of your report).

\section{Results}
\label{sec:result}

We hypothesize that as the distance from a substance abuse treatment facility increases, the number of 911 calls also increases. This means that the further someone is from a treatment center, the less access they have to treat drug addiction, so the more likely they would exhibit drug-induced behavior that results in someone calling 911. Our null hypothesis is that there is no relation between distance to an SATC and the likelihood of calling 911. 

We conducted a set paired t-test to examine the statistical likelihood that one distance ring has a greater mean number of calls than the next larger ring. The pairwise comparison table is seen in Table~\ref{tabl:Table1}.

\begin{table}[h]
\scriptsize
\begin{center}
\documentclass{article}
\usepackage{multirow}
\usepackage{amsmath}
\usepackage{ulem}
\usepackage[table]{xcolor}
\begin{document}
\begin{table}[!h]
\centering
\begin{tabular}{lllllll}
\cline{1-7}
\multicolumn{1}{c}{} &
  \multicolumn{1}{|r}{Mean\_1} &
  \multicolumn{1}{r}{Mean\_2} &
  \multicolumn{1}{r}{Difference} &
  \multicolumn{1}{r}{p\_value} &
  \multicolumn{1}{r}{Tailed\_p} &
  \multicolumn{1}{r}{stars} \\
\cline{1-7}
\multicolumn{1}{l}{ttest dist\_group\_100 == dist\_group\_250} &
  \multicolumn{1}{|r}{2838.745} &
  \multicolumn{1}{r}{1870.036} &
  \multicolumn{1}{r}{968.709} &
  \multicolumn{1}{r}{0.056} &
  \multicolumn{1}{r}{0.028} &
  \multicolumn{1}{r}{*} \\
\multicolumn{1}{l}{ttest dist\_group\_2250 == dist\_group\_2500} &
  \multicolumn{1}{|r}{206.033} &
  \multicolumn{1}{r}{135.579} &
  \multicolumn{1}{r}{70.455} &
  \multicolumn{1}{r}{0.004} &
  \multicolumn{1}{r}{0.002} &
  \multicolumn{1}{r}{***} \\
\multicolumn{1}{l}{ttest dist\_group\_250 == dist\_group\_500} &
  \multicolumn{1}{|r}{1870.036} &
  \multicolumn{1}{r}{1680.406} &
  \multicolumn{1}{r}{189.630} &
  \multicolumn{1}{r}{0.491} &
  \multicolumn{1}{r}{0.245} &
  \multicolumn{1}{r}{} \\
\multicolumn{1}{l}{ttest dist\_group\_500 == dist\_group\_750} &
  \multicolumn{1}{|r}{1680.406} &
  \multicolumn{1}{r}{1069.475} &
  \multicolumn{1}{r}{610.931} &
  \multicolumn{1}{r}{0.007} &
  \multicolumn{1}{r}{0.003} &
  \multicolumn{1}{r}{***} \\
\multicolumn{1}{l}{ttest dist\_group\_750 == dist\_group\_1000} &
  \multicolumn{1}{|r}{1069.475} &
  \multicolumn{1}{r}{979.865} &
  \multicolumn{1}{r}{89.610} &
  \multicolumn{1}{r}{0.392} &
  \multicolumn{1}{r}{0.196} &
  \multicolumn{1}{r}{} \\
\multicolumn{1}{l}{ttest dist\_group\_1000 == dist\_group\_1250} &
  \multicolumn{1}{|r}{979.865} &
  \multicolumn{1}{r}{823.863} &
  \multicolumn{1}{r}{156.002} &
  \multicolumn{1}{r}{0.006} &
  \multicolumn{1}{r}{0.003} &
  \multicolumn{1}{r}{***} \\
\multicolumn{1}{l}{ttest dist\_group\_1250 == dist\_group\_1500} &
  \multicolumn{1}{|r}{823.863} &
  \multicolumn{1}{r}{563.835} &
  \multicolumn{1}{r}{260.028} &
  \multicolumn{1}{r}{0.000} &
  \multicolumn{1}{r}{0.000} &
  \multicolumn{1}{r}{***} \\
\multicolumn{1}{l}{ttest dist\_group\_1500 == dist\_group\_1750} &
  \multicolumn{1}{|r}{563.835} &
  \multicolumn{1}{r}{412.993} &
  \multicolumn{1}{r}{150.842} &
  \multicolumn{1}{r}{0.040} &
  \multicolumn{1}{r}{0.020} &
  \multicolumn{1}{r}{**} \\
\multicolumn{1}{l}{ttest dist\_group\_1750 == dist\_group\_2000} &
  \multicolumn{1}{|r}{412.993} &
  \multicolumn{1}{r}{248.251} &
  \multicolumn{1}{r}{164.742} &
  \multicolumn{1}{r}{0.001} &
  \multicolumn{1}{r}{0.001} &
  \multicolumn{1}{r}{***} \\
\multicolumn{1}{l}{ttest dist\_group\_2000 == dist\_group\_2250} &
  \multicolumn{1}{|r}{248.251} &
  \multicolumn{1}{r}{206.033} &
  \multicolumn{1}{r}{42.218} &
  \multicolumn{1}{r}{0.039} &
  \multicolumn{1}{r}{0.020} &
  \multicolumn{1}{r}{**} \\
\cline{1-7}
\end{tabular}

\footnotesize{
*** p$<$.01, ** p$<$.05, * p$<$.1
}
\end{table}
\end{document}

\end{center}
\flushleft
\caption{\textbf{Two Sample T-Tests}}
\label{tab:Table1}
\end{table}

The results support the rejection of the null hypothesis for seven of the paired tests. Five of the pair t-tests exhibit a p-value less than 0.01, indicating statistical significance between sequential distance rings. In addition, two of the tests show p-values less than 0.05. This is still statistically significant, although less strong than the groups mentioned above. 

For three of the paired t-test, we cannot reject the null hypothesis as the null hypothesis value falls within the 95\% confidence interval. The 100-250 meter ring test, has a p-value of 0.056, which is close to the threshold of significance. This suggests that there may be a weak effect presence, so further tests should be done.

However, the 250-500 meter and 750-1000 meter ring tests exhibit p-values of 0.491 and 0.392 respectively. These high p-values show no evidence of statistical significance in the mean values of the distance rings, likely attributing the differences in the mean values to random variation.

To account for SATC within the unit of analysis, we analyze the mean calls per distance group across SATC. The results are in Figure~\ref{fig:Figure1}. Figure~\ref{fig:Figure1} is overlapped with a confidence interval. Although the figure mirrors Figure~\ref{fig:Figure1} visually, the majority of the confidence intervals overlap some range. 
\begin{figure}{h!}
    \centering
\includegraphics[width=0.75\linewidth]{CI_Graph.png}
    \caption{Mean Calls per SATC per distance group}
    \label{fig:Figure1}
     \textit{X-axis consist of the different distance groups used. Distance measured in meters.}
    \textit{Y-axis is measured in mean calls per kilometers squared.}
\end{figure}
The confidence intervals of 250-500 meters and 750-1000 meters significantly overlap. The other confidence intervals do overlap with the sequential intervals, but not visually as significantly as in the two comparisons mentioned above. These observations additionally support the p-values observed in Table 1. 

\begin{figure}{h!}
    \centering
\includegraphics[width=0.75\linewidth]{Crim_CI_Graph.png}
    \caption{Mean Crime Calls per SATC per distance group}
    \label{fig:Figure2}
\end{figure}
%\subsubsection{Summary Statistics}
%State your working hypothesis and null hypothesis - Explain why the statistics you’re presenting are appropriate -Present statistics and explain what they mean (e.g., do your results support the working hypothesis? 

%Explain what analyses you did, provide evidence (like in the descriptive stats exercise, but refined and clear) and then explain what your results mean.


\section{Discussion}
\label{sec:discussion}

One explanation for this relationship could be due to the positive feedback loop between drugs, crime, and lower standards of living in those places. SATC may be put into neighborhoods where there are a lot of drug users so it is accessible and convenient to those who need it. Additionally, neighborhoods with higher rates of drug usage tend to have a lower standard of living. Another indication of a lower standard of living is higher crime rates, which could potentially contribute to higher counts of 911 calls closer to an SATC. 

Additionally, SATC may be strategically placed in lower standard-of-living neighborhoods due to pressure from higher standard-of-living communities. Richer communities tend to exhibit a "Not In My Backyard" (NIMMBY) mentality, opposing the placement of treatment centers in their neighborhoods. This resistance stems from the fear or stigma residents have about drug users. They might perceive SATC as bringing dangerous or mentally ill people into their area, impacting the safety of the area. Thus, higher standard-of-living neighborhoods tend to lobby for NIMBY policies in their communities, pushing SATC to be placed in areas that can't or don't lobby these same interests. Additionally, higher standard-of-living areas tend to have lower crime rates in general. This creates a positive feedback system in these neighborhoods, potentially explaining why the farther a distance ring is from a SATC, the fewer mean 911 calls there are. Those distance rings could be on the cusp or or geographically in neighborhoods that don't have SATC. 
    
Another explanation for the lower frequency of 911 calls in further distance rings could be due to our data not counting 911 calls twice. When we geo-coded the 911 calls, we assigned each observation to only one SATC, the one geographically closest to it. This was to limit the amount of overlap that some SATC might have at the smaller distance rings, as it could exaggerate our data frequencies. However, in farther distance rings, a call might be similar in distance to two SATCs, but will only be assigned to be correlated to one, when the reality it could be correlated to either. This thins the observations and data associated in farther-out distance rings. 

One significant limitation was the inability to fine-tune 911 call data to determine whether calls were directly related to drug activity. Without additional data, such as toxicology reports or detailed incident descriptions, it remains challenging to assess the role of drugs in criminal behavior or whether the crimes involved were violent in nature. A better understanding of what calls were drug-related or induced could provide more accurate results on the relationship between drug abuse and crime, and the role of SATC in amplifying or mitigating this. 
    
Another limitation is the population density of this urban sample. The observed pattern of 911 calls could have been influenced by higher population density in more urban parts of Detroit, which naturally results in an increase of emergency calls. Future research should explore similar analyses in areas with varying population densities to better control for this variable and isolate the effects of proximity to substance abuse treatment centers. 

Another limitation was the inability to create a regression line (or line of best fit) for this relationship. Since we collapsed the data multiple times, we lost variability to perform a regression line. First, we collapsed the data by count to get the total counts for each distance group. Then we collapsed the data by grouping by SATC designation, to get the total observations per distance group per SATC. Finally, we collapsed that data set by mean, creating mean counts of 911 calls per distance group adjusted for SATC correlation. This left one observation row that created Figure 4. We lost analysis flexibility due to this data manipulation approach. In the future, a different data manipulation approach might leave more room for further data analysis. 

We analyzed data from 2017 only. This leaves room for error from using a small dataset. We see future research supporting these findings or supporting a stronger pattern through these two ways: Firstly, performing cross-sectional data within multiple different years. Then, comparing the mean calls per SATC per distance group and identifying trend lines. This option could strengthen the correlational relationships we found in this paper. Secondly, a paper could analyze time series based on the openings and closings of substance abuse treatment centers. This option might identify a causal impact of a SATC's presence might have on a neighborhood as it pertains to 911 calls, and subsequently, crime rates. These results could impact policy-making decisions when it comes to where to fund and support the placement of SATC to serve populations in need. 

\section{Conclusion}
\label{sec:conclusion}

    We hypothesized that as the distance from a SATC increases, the frequency of 911 calls also increases. The closer a SATC is, the more accessible treatment centers are, which should reduce crime rates in the area, measured by 911 calls. We conducted a cross-sectional analysis to examine the overlap of 911 calls' geographical locations to their distances to their closest SATC. Our results show that contrary to our hypothesis, as distance from a SATC increases, the amount of 911 calls decreases. This distance decay relationship indicates that there is a relationship between SATC placements and 911 calls, but just in the opposite direction of the original hypothesis. 
   % A


 
%Re-state (in different words) what you did and what you learned. If your discussion (Section 6) would be short, you can just have a Conclusion section that includes your discussion (that is, leave out a separate Discussion section).

\newpage
\singlespacing
\setlength\bibsep{0pt}

\bibliographystyle{chicago}
\bibliography{reference.bib}
https://www.pewresearch.org/short-reads/2024/04/24/what-the-data-says-about-crime-in-the-us/

\newpage
\section*{Data Appendix} \label{sec:appendixa}
\addcontentsline{toc}{section}{Appendix A}

Our replication package is linked \href{https://github.com/ecn310/course-project-zipcentercrime/tree/main/Reproducibility%20Package}{here}.

We isolated our data from both Police Serviced 911 Calls and N-SSATS to only include observations from the year 2017. We picked this year due to an overlap in data and no potential disruptions from COVID-19. The year 2018 or 2019 from both data sets could additionally work. 

We used the ArcGISPro application to layer both 2017 datasets and geocode their locations. Using the coordinates (longitude and latitude variables) of each point, we created two new variables. Every 911 call is assigned one SATC, the one geographically closest to it. This means no observation is counted twice. 

We used Stata 18.0 to identify distance buckets, referred to as distance groups, and perform our analysis. We normalize the data to account for how differences in total area of a distance group might alter results. We did this by dividing the area by the total count of 911 calls in each distance group. 


% \begin{table}[h]
%\centering
%\caption{Ratio Analysis}
%begin{tabular}{l||c|c|cc|c|c}
%\hline\hline
 %           &ratio\_results&            &            &            &            &            \\
     %       &       Ratio&   Std. Err.&  [95 5Conf.&   Interval]&     T score&     P value\\
%\hline
%100/250     &    1.518016&     .305492&     .901933&      2.1341&    4.969087&    .0000112\\
%250/500     &    1.112848&    .1703623&    .7692794&    1.456416&    6.532241&    6.14e-08\\
%500/750     &    1.571244&    .1966827&    1.174596&    1.967893&    7.988725&    4.92e-10\\
%750/1000    &    1.091451&    .1104992&    .8686081&    1.314294&    9.877456&    1.25e-12\\
%1000/1250   &    1.189354&    .0748615&    1.038382&    1.340327&     15.8874&    1.32e-19\\
%1250/1500   &    1.461179&    .1268808&    1.205299&    1.717058&    11.51615&    1.00e-14\\
%1500/1750   &    1.365242&    .1935897&    .9748308&    1.755652&    7.052244&    1.08e-08\\
%1750/2000   &     1.66361&    .2204068&    1.219117&    2.108103&    7.547907&    2.09e-09\\
%2000/2250   &    1.204906&     .110774&    .9815093&    1.428303&    10.87716&    6.33e-14\\
%2250/2500   &    1.519659&    .1688579&    1.179124&    1.860193&     8.99963&    1.91e-11\\
%\hline\hline
%\end{tabular}
%\flushleft
%\begin{footnotesize}
%\begin{singlespace}
%\textbf{Note 1}: A larger T score indicates the estimated coefficient is more significantly different from the hypothesized value.
%\textbf{Note 2}: A smaller P value p<0.05 means the result is statistically significant. This means the estimated coefficient is likely not zero.
%\end{singlespace}
%\end{footnotesize}
%\end{table} %

\end{document}